
\section{Causal Analysis of Hidden State Monitoring}

\subsection{Causal Model}

We formalize the hidden state monitoring mechanism using Pearl's structural causal model framework \cite{pearl2009causality}. The causal graph is defined as:

\begin{equation}
X \rightarrow H \rightarrow Y
\end{equation}

where $X$ represents the input (harmful/benign), $H$ represents the hidden state, and $Y$ represents the output (safe/harmful).

\subsection{Causal Intervention Analysis}

Using do-calculus, we compute the causal effect of intervention:

\begin{equation}
P(Y=\text{harmful} \mid \text{do}(R < \tau)) = 0.000
\end{equation}

The Average Causal Effect (ACE) is:
\begin{equation}
\text{ACE} = E[Y \mid \text{do}(X=1)] - E[Y \mid \text{do}(X=0)] = 0.734
\end{equation}

\subsection{Mediation Analysis}

We decompose the total effect into direct and indirect effects:

\begin{itemize}
    \item Natural Direct Effect (NDE): 0.299
    \item Natural Indirect Effect (NIE): 0.458
    \item Proportion Mediated: 60.5%
\end{itemize}

The high proportion mediated (60.5%) confirms that the harmful effect is primarily transmitted through hidden states, validating our monitoring approach.

\subsection{Counterfactual Analysis}

The Effect of Treatment on the Treated (ETT) measures the benefit of intervention for harmful inputs:

\begin{equation}
\text{ETT} = E[Y(1) - Y(0) \mid X=1] = 0.867
\end{equation}

This shows that 86.7% of harmful inputs would be successfully blocked by the intervention.
